\Needspace{42\baselineskip}
\section{Groupoïde de l'extension bidirectionnelle et algèbre opératorielle}
\label{sec:cpe-grupoide-two-way}

Soit \(X_F^{\leftrightarrow}\) l'extension naturelle de l'état enrichi et
\(\sigma\) son homéomorphisme de décalage. On définit le groupoïde de
transformation
\begin{equation}
 \mathcal G_{\rm TPK}
 :=X_F^{\leftrightarrow}\rtimes_\sigma\mathbb Z,
 \label{eq:cpe-grupoide-tpk}
\end{equation}
avec
\[
 s(x,n)=x,
 \quad r(x,n)=\sigma^n x,
 \quad (\sigma^nx,m)(x,n)=(x,m+n),
 \quad (x,n)^{-1}=(\sigma^nx,-n).
\]

\begin{theorem}[Complétion groupoïdale du transport bidirectionnel]
\label{thm:cpe-grupoide-tpk}
Le groupoïde \(\mathcal G_{\rm TPK}\) est localement compact, séparé,
à base dénombrable, étale et moyennable. Pour la mesure projective invariante et
non atomique \(\mu_{\rm HMT}\), si le décalage est ergodique et
essentiellement libre,
\[
 \mathcal M_{\rm TPK}
 :=L^\infty(X_F^{\leftrightarrow},\mu_{\rm HMT})
 \rtimes_\sigma\mathbb Z
\]
est le facteur hyperfini de type \(II_1\).
\end{theorem}

\begin{proof}
L'action d'un groupe discret dénombrable par homéomorphismes sur un compact
métrisable produit un groupoïde étale, localement compact, séparé et
à base dénombrable. La moyennabilité de \(\mathbb Z\) implique la moyennabilité
du groupoïde. La mesure projective est invariante par le décalage.
L'ergodicité et la liberté essentielle rendent le produit croisé factoriel ; la non-
atomicité exclut le cas matriciel de type I. La trace induite par
\(\mu_{\rm HMT}\) est finie et fidèle, et la moyennabilité produit l'
hyperfinitude. Par conséquent, le facteur est de type \(II_1\).
\end{proof}

La classification modulaire de type \(III\) appartient à une mesure
quasi invariante distincte, munie d'un cocycle de Radon--Nikodym déclaré. Elle ne se
déduit pas de la simple présence d'une dynamique KMS et ne rouvre pas la complétion
groupoïdale précédente.
